\documentclass[conference]{IEEEtran}
\IEEEoverridecommandlockouts

\usepackage{cite}
\usepackage{amsmath,amssymb,amsfonts}
\usepackage{algorithmic}
\usepackage{graphicx}
\usepackage{textcomp}
\usepackage{xcolor}
\usepackage{booktabs}
\usepackage{multirow}
\usepackage{array}
\usepackage{url}
\usepackage{hyperref}
\usepackage{microtype}

\hyphenation{pre-dic-tive main-te-nance im-bal-ance XG-Boost Light-GBM SMOTE LSTM}

\def\BibTeX{{\rm B\kern-.05em{\sc i\kern-.025em b}\kern-.08em
    T\kern-.1667em\lower.7ex\hbox{E}\kern-.125emX}}

\begin{document}

\title{Hybrid Machine Learning and Deep Learning for Health Status
       Classification and Remaining Useful Life Prediction in Manufacturing}

\author{
Achmad Zikran Maulida\textsuperscript{1},
Amir Hamzah\textsuperscript{2},
Aqsa Zamzami\textsuperscript{3},
Reynaldi Chandra Kirana\textsuperscript{4} \\

\textsuperscript{}Jurusan Teknik Informatika dan Komputer,
Program Studi Teknik Informatika \\
Politeknik Negeri Jakarta \\
Jln. Prof.\ DR.\ G.A.\ Siwabessy, Kampus Universitas Indonesia, Depok 16425 \\

\textsuperscript{1}achmad.zikran.maulida.tik23@stu.pnj.ac.id \quad
\textsuperscript{2}amir.hamzah.tik23@stu.pnj.ac.id \\
\textsuperscript{3}aqsa.zamzami.tik23@stu.pnj.ac.id \quad
\textsuperscript{4}reynaldi.chandra.kirana.tik23@stu.pnj.ac.id
}
\maketitle

\begin{abstract}
Machine failure events constitute as few as 0.056\% of industrial
IoT sensor observations, creating an extreme class imbalance that,
when compounded by temporal data leakage from naive train-test
splitting, renders most standard predictive maintenance pipelines
unreliable under real deployment conditions. This paper presents
PRIME (Predictive Reliability and Intelligence Maintenance Engine),
an end-to-end pipeline that resolves both challenges through three
principled contributions: (1)~\textbf{Stratified Sequential Block
Sampling (SSBS)}, a temporally-aware undersampling strategy that
reduces HEALTHY-class dominance from 97.4\% to 87.1\% while
preserving chronological data integrity; (2)~\textbf{Temporal Label
Engineering} with a P90-based Sensor Confirmation Layer that
converts binary failure signals into three operationally actionable
health states (HEALTHY, WARNING, CRITICAL); and (3)~a
\textbf{cascaded XGBoost--LSTM inference architecture} that
activates the regression module conditionally upon anomaly
detection, reducing real-time inference overhead by
$\sim$70\% compared to always-on sequential execution. Evaluated
on a 100,000-record, 20-machine time-series dataset under strict
machine-based splitting with train-only SMOTE augmentation, the
pipeline achieves a WARNING-class F1-score of 0.9818 with
\textbf{zero fatal misclassifications} (CRITICAL$\rightarrow$HEALTHY),
and predicts Remaining Useful Life with an MAE of 0.7985~days,
with 98.04\% of predictions falling within one day of ground truth.
These results establish that principled temporal handling of
severely imbalanced sensor data enables production-grade predictive
maintenance without sacrificing computational efficiency.
\end{abstract}

\begin{IEEEkeywords}
Predictive Maintenance, XGBoost, LSTM, Temporal Label Engineering, SSBS,
Remaining Useful Life, Time-Series Classification, SMOTE
\end{IEEEkeywords}

% ============================================================
\section{INTRODUCTION}
% ============================================================

The emergence of Industry~4.0 has enabled pervasive deployment of
Internet-of-Things (IoT) sensors across manufacturing facilities, generating
continuous high-frequency streams of machine operational data
\cite{lee2014industryiot}. Traditional maintenance paradigms---corrective
(fix on failure) and preventive (fixed-schedule)---are increasingly
inadequate: the former incurs costly unplanned downtime, while the latter
wastes resources on unnecessary interventions. Predictive Maintenance (PdM)
leverages sensor data and machine learning to schedule maintenance precisely
when needed, transforming reactive operations into data-driven,
condition-based workflows \cite{zhang2019pdmsurvey}.

Real-world PdM deployments face two compounding technical challenges. First,
failure events are exceedingly rare: in the dataset analyzed in this work,
only 56 records out of 100,000 total observations carry a failure label---a
rate of 0.056\%. This extreme class imbalance causes standard classifiers to
be overwhelmingly biased toward predicting the majority HEALTHY class
\cite{chawla2002smote}. Resampling via SMOTE \cite{chawla2002smote} is
commonly applied, but if executed before train-test splitting it introduces
temporal data leakage \cite{kaufman2012leakage}. Second, the temporal
structure of time-series data renders standard cross-validation invalid:
conventional $k$-fold validation violates temporal causality by including
future observations in training folds, producing inflated performance
estimates \cite{bergmeir2018tscv}.

Existing literature lacks a unified framework that simultaneously addresses
extreme class imbalance without temporal leakage, multi-class label
engineering from binary failure signals, and a cascaded ML+DL inference
architecture that is computationally efficient. This paper fills this gap with
the following contributions:

\begin{enumerate}
\item \textbf{Stratified Sequential Block Sampling (SSBS)}: a
  temporally-aware undersampling strategy that selects chronological blocks
  per machine, reducing HEALTHY dominance from 97.36\% to 87.09\% without
  destroying time-series causality.
\item \textbf{Three-class Temporal Label Engineering} with an empirically
  calibrated Sensor Confirmation Layer (P90 thresholds, minimum 2 sensors),
  transforming binary failure labels into actionable
  HEALTHY/WARNING/CRITICAL classifications.
\item \textbf{Machine-based splitting combined with train-only SMOTE}:
  an enterprise-grade zero-leakage pipeline for time-series PdM modeling.
\item \textbf{Cascaded XGBoost+LSTM deployment architecture} achieving
  state-of-the-art classification (F1$=0.9906$) and RUL prediction
  (MAE$=0.7985$~days, Error$\leq$1d$=98.04\%$) while reducing inference
  overhead by $\sim$70\% through conditional LSTM activation.
\end{enumerate}

% ============================================================
\section{RELATED WORK}
% ============================================================

\subsection{Deep Learning for RUL Prediction}
Long Short-Term Memory (LSTM) networks have become the dominant architecture
for RUL prediction from sequential sensor data. Hochreiter and
Schmidhuber's foundational LSTM \cite{hochreiter1997lstm} addressed the
vanishing gradient problem in vanilla RNNs, enabling learning of long-range
temporal dependencies critical for degradation modeling. Zheng et~al.\
\cite{zheng2017lstm_rul} demonstrated LSTM superiority over feed-forward
networks on the CMAPSS turbofan dataset. Hybrid CNN-LSTM architectures
\cite{li2018cnnlstm} extract local features (CNN) before temporal modeling
(LSTM), outperforming pure LSTM on multi-sensor industrial datasets. Gated
Recurrent Units (GRU) \cite{cho2014gru} offer a computationally lighter
alternative, but empirical studies on complex multi-sensor degradation
signals consistently favor LSTM for asymmetric temporal patterns.

\subsection{Tree-Based Methods for Fault Classification}
Ensemble tree methods remain state-of-the-art for tabular and time-series
classification in industrial fault detection. Chen and Guestrin
\cite{chen2016xgboost} introduced XGBoost, a scalable gradient-boosted
framework consistently outperforming alternatives on structured data. Random
Forest \cite{breiman2001rf}, while a strong baseline, tends to produce wider
decision boundaries that inflate false alarm rates in imbalanced datasets.
LightGBM \cite{ke2017lightgbm} achieves faster training via leaf-wise growth,
but its aggressive early stopping causes premature convergence when minority
class samples are sparse, as observed in our experiments.

\subsection{Imbalance Handling and Splitting Strategies}
Class imbalance is ubiquitous in PdM datasets \cite{li2020pdmimbalance}.
SMOTE \cite{chawla2002smote} generates synthetic minority samples by feature-
space interpolation, but applying SMOTE before splitting introduces leakage:
synthetic samples may encode statistical properties of future test data,
producing overly optimistic evaluation results \cite{blagus2013smote_bias}.
For multi-entity time-series, machine-based splitting---assigning complete
machine histories to distinct partitions---is more appropriate than global
temporal splits, as it evaluates generalization across unseen machines
\cite{bergmeir2018tscv}. The temporal leakage risk from early SMOTE
application and the efficiency of conditional LSTM activation have not been
systematically studied in a unified PdM pipeline, constituting the research
gap addressed in this work.

% ============================================================
\section{RESEARCH METHODOLOGY}
% ============================================================

\subsection{Dataset Description}

The experimental dataset (\texttt{sensor\_readings.csv}) comprises continuous
hourly IoT sensor readings from 20 identical industrial machines over 208~days
(1~July~2025 to 25~January~2026). Each machine contributes exactly 5,000
timestamped records, yielding 100,000 total observations. The dataset exhibits
extreme class imbalance: only 56~records are labeled \texttt{failure=1}
(0.056\%), with an average of 2.80 failure events per machine. Eight sensor
features span thermal, mechanical, electrical, and environmental domains
(Table~\ref{tab:dataset_stats}). A supplementary
\texttt{maintenance\_logs.csv} (500~records) provides technician notes used
for NLP-derived feature enrichment.

\begin{table}[htbp]
\caption{TABLE I: Dataset Statistics}
\centering
\begin{tabular}{ll}
\hline
\textbf{Property} & \textbf{Value} \\
\hline
Total Records          & 100,000 \\
Number of Machines     & 20 \\
Observation Period     & 208 days \\
Records per Machine    & 5,000 \\
Sensor Features        & 8 \\
Failure Events         & 56 (0.056\%) \\
Healthy Events         & 99,944 (99.944\%) \\
Temporal Resolution    & 1 hour \\
\hline
\end{tabular}
\label{tab:dataset_stats}
\end{table}

\subsection{Temporal Label Engineering and Sensor Confirmation}

The binary \texttt{failure} label is insufficient for operational
decision-making. We transform it into a three-class ordinal label using a
temporal backward-labeling algorithm: for each failure event at timestamp
$T_f$, records in $[T_f - W_{\text{WARNING}},\,T_f)$ are labeled WARNING and
records in $[T_f - W_{\text{CRITICAL}},\,T_f)$ are relabeled CRITICAL, where
$W_{\text{WARNING}} = 48$~h and $W_{\text{CRITICAL}} = 24$~h. These windows
were derived empirically from Failure Autopsy analysis of machines M-01,
M-05, and M-09.

A \textbf{Sensor Confirmation Layer} further refines the WARNING label: a
WARNING is retained only if at least 2 of the 6 high-priority sensors exceed
their respective P90 thresholds computed from the HEALTHY distribution
(Table~\ref{tab:p90}). This step downgraded 49 WARNING rows to HEALTHY
(1.82\%), while all CRITICAL rows were retained.

\begin{equation}
y_i = \begin{cases}
2 & \text{if } T_f - W_{\text{CRIT}} \leq t_i < T_f \\
1 & \text{if } T_f - W_{\text{WARN}} \leq t_i < T_f - W_{\text{CRIT}} \\
  & \quad \text{and } \sum_{s=1}^{6} \mathbf{1}[x_{i,s} > \theta_s^{P90}] \geq 2 \\
0 & \text{otherwise}
\end{cases}
\label{eq:label}
\end{equation}

%%====================================================================
%% FIGURE PLACEHOLDER B: Sensor Timeline Menuju Failure
%% Sumber: notebooks/fase_2_eda/02_eda_forensik.ipynb
%% Export sebagai: figures/sensor_timeline_failure.png
%% Di Overleaf: upload ke folder "figures/", uncomment blok berikut.
%%====================================================================
% \begin{figure}[htbp]
%   \centering
%   \includegraphics[width=\columnwidth]{figures/sensor_timeline_failure.png}
%   \caption{Sensor trajectory for machine M-05 during the 72~hours
%            preceding failure. Degradation onset visible at $T{-}48$h
%            (WARNING boundary) with rapid escalation at $T{-}24$h.}
%   \label{fig:sensor_timeline}
% \end{figure}
%%====================================================================

\begin{table}[htbp]
\caption{TABLE II: Final Label Distribution After Temporal Engineering}
\centering
\begin{tabular}{lrr}
\hline
\textbf{Class} & \textbf{Count} & \textbf{Percentage} \\
\hline
HEALTHY (0)  & 97,364 & 97.364\% \\
WARNING (1)  &  1,247 &  1.247\% \\
CRITICAL (2) &  1,389 &  1.389\% \\
\hline
\textbf{Total} & \textbf{100,000} & \textbf{100.000\%} \\
\hline
\end{tabular}
\label{tab:label_dist}
\end{table}

\begin{table}[htbp]
\caption{TABLE III: P90 Sensor Thresholds --- Confirmation Layer}
\centering
\begin{tabular}{ll}
\hline
\textbf{Sensor} & \textbf{P90 Threshold (HEALTHY)} \\
\hline
temperature        & 76.30 $^\circ$C \\
vibration          & 0.59 mm/s \\
pressure           & 103.80 bar \\
rpm                & 2,540.00 \\
power\_consumption & 82.10 kW \\
noise\_level       & 74.30 dB \\
\hline
\end{tabular}
\label{tab:p90}
\end{table}

All six primary sensors exhibit large Cohen's $d$ effect sizes ($d > 2.0$),
confirming high discriminative power. Vibration ranks highest ($d=3.37$),
followed by pressure ($d=3.06$), rpm ($d=2.91$), noise\_level ($d=2.89$),
temperature ($d=2.71$), and power\_consumption ($d=2.61$). Operating hours
($d=0.29$) and humidity ($d=0.21$) show medium effect and are retained but
deprioritized.

\subsection{Feature Engineering}

Raw sensor values alone provide limited temporal context. We enrich the
feature space through five systematic groups, yielding 69 final features
(Table~\ref{tab:features}).

\begin{table}[htbp]
\caption{TABLE IV: Feature Engineering Summary (69 Input Features)}
\centering
\begin{tabular}{lll}
\hline
\textbf{Group} & \textbf{Count} & \textbf{Description} \\
\hline
Original Sensors    &  8 & Raw sensor values \\
Rolling Statistics  & 36 & 6 sensors $\times$ 2 windows $\times$ 3 stats \\
Lag Features        & 18 & 6 sensors $\times$ lag \{1, 6, 24\}h \\
Cross-Sensor Ratios &  4 & Pairwise sensor interaction \\
Degradation Proxy   &  1 & hours\_since\_last\_maintenance \\
NLP-Derived         &  2 & damage\_category, severity\_score \\
\hline
\textbf{TOTAL} & \textbf{69} & After label/ID columns removed \\
\hline
\end{tabular}
\label{tab:features}
\end{table}

Rolling Statistics (window sizes 24h and 48h, computing mean, standard
deviation, and maximum for each of the six high-priority sensors) constitute
the top-ranked features by XGBoost importance. All 69 numeric features are
standardized using \texttt{StandardScaler} fit exclusively on training
machines (M-01 to M-14).

\subsection{Imbalance Handling: SSBS and SMOTE}

\textbf{Stratified Sequential Block Sampling (SSBS).}
Naive random undersampling destroys temporal continuity, breaking the
sequential structure required by lag and rolling features. SSBS selects two
chronological blocks per machine: \textbf{Block~A} (first 500 rows, capturing
prime operating conditions) and \textbf{Block~B} (last 500 rows with a 24-h
safety buffer before the first WARNING zone). All WARNING and CRITICAL rows
are always retained in full. The SSBS formula per machine $m$ is:
\begin{equation}
S_m^{\text{HEALTHY}} = \text{Block}_A(m,500) \cup \text{Block}_B(m,500,\delta=24\text{h})
\label{eq:ssbs}
\end{equation}

%%====================================================================
%% FIGURE PLACEHOLDER C: SSBS Block Selection Diagram
%% Sumber: notebooks/fase_6_imbalance/06_imbalance_handling.ipynb
%% Export sebagai: figures/ssbs_diagram.png
%%====================================================================
% \begin{figure}[htbp]
%   \centering
%   \includegraphics[width=\columnwidth]{figures/ssbs_diagram.png}
%   \caption{SSBS block selection per machine. Block~A captures prime
%            operating conditions and Block~B captures the pre-failure
%            boundary with a 24-h safety buffer.}
%   \label{fig:ssbs_diagram}
% \end{figure}
%%====================================================================

\begin{table}[htbp]
\caption{TABLE V: SSBS Result Distribution}
\centering
\begin{tabular}{lrrl}
\hline
\textbf{Class} & \textbf{Before} & \textbf{After} & \textbf{After \%} \\
\hline
HEALTHY  & 97,364 & 17,787 & 87.09\% \\
WARNING  &  1,247 &  1,247 &  6.11\% \\
CRITICAL &  1,389 &  1,389 &  6.80\% \\
\hline
\textbf{Total} & \textbf{100,000} & \textbf{20,423} & \textbf{100\%} \\
\hline
\end{tabular}
\label{tab:ssbs}
\end{table}

\textbf{Machine-Based Split.} Complete machine histories are assigned to
each partition to guarantee that every split contains full run-to-failure
cycles and that the model generalizes across unseen machines:

\begin{table}[htbp]
\caption{TABLE VI: Machine-Based Dataset Split}
\centering
\begin{tabular}{llrl}
\hline
\textbf{Split} & \textbf{Machines} & \textbf{Samples} & \textbf{\%} \\
\hline
Train      & M-01 to M-14      & 14,419 & 70.60\% \\
Validation & M-15, M-16, M-17  &  3,288 & 16.10\% \\
Test       & M-18, M-19, M-20  &  2,716 & 13.30\% \\
\hline
\textbf{Total} & 20 machines & \textbf{20,423} & \textbf{100\%} \\
\hline
\end{tabular}
\label{tab:split}
\end{table}

\textbf{Train-Only SMOTE.} Residual class imbalance in the training set
(HEALTHY: 12,504; WARNING: 901; CRITICAL: 1,014) is addressed by SMOTE
\cite{chawla2002smote} applied \textit{exclusively} to the training partition
after splitting, targeting WARNING and CRITICAL to 4,000 samples each. The
6,085 synthetic samples generated had zero influence on the 3,288 validation
and 2,716 test records, ensuring unbiased evaluation.

\begin{table}[htbp]
\caption{TABLE VII: SMOTE Effect on Training Set}
\centering
\begin{tabular}{lrrr}
\hline
\textbf{Class} & \textbf{Before} & \textbf{After} & \textbf{Synthetic} \\
\hline
HEALTHY  & 12,504 & 12,504 &     0 \\
WARNING  &    901 &  4,000 & 3,099 \\
CRITICAL &  1,014 &  4,000 & 2,986 \\
\hline
\textbf{Total} & \textbf{14,419} & \textbf{20,504} & \textbf{6,085} \\
\hline
\end{tabular}
\label{tab:smote}
\end{table}

\subsection{Predictive Models}

\textbf{Model 1 --- Health Status Classifier.}
Three candidate algorithms were trained on the SMOTE-augmented set
(20,504~samples, 69~features) and assessed on natural validation and test
sets. Hyperparameters are summarized in Table~\ref{tab:hyper_clf}.

\begin{itemize}
\item \textit{Random Forest} (Baseline): 300 trees,
  \texttt{class\_weight='balanced'}, F1$_\text{val}=0.9292$, WARNING
  F1$_\text{val}=0.811$. High variance gap indicates excessive false alarms.

\item \textit{XGBoost V2} (Selected): \texttt{n\_estimators=1000},
  \texttt{max\_depth=4}, \texttt{learning\_rate=0.01},
  \texttt{subsample=0.8}, \texttt{reg\_alpha=0.5}, \texttt{reg\_lambda=2.0},
  converging at iteration 499. A WARNING probability threshold of 0.60
  (vs.\ default 0.50) eliminates 953 of 954 false alarms while maintaining
  WARNING F1$=0.9818$ on validation.

\item \textit{LightGBM} (Baseline): converged aggressively at iteration 27
  due to logloss stagnation under sparse minority samples. Competitive after
  threshold tuning (0.65) but inferior on the primary validation WARNING F1
  metric.
\end{itemize}

\begin{table}[htbp]
\caption{TABLE VIII: Model 1 Hyperparameter Summary}
\centering
\begin{tabular}{p{2.2cm}ccc}
\hline
\textbf{Parameter} & \textbf{Rand.~Forest} & \textbf{XGBoost V2} & \textbf{LightGBM} \\
\hline
n\_estimators   & 300      & 499$^\dagger$ & 27$^\dagger$ \\
max\_depth      & 20       & 4             & N/A \\
learning\_rate  & N/A      & 0.01          & 0.05 \\
subsample       & N/A      & 0.8           & N/A \\
reg\_alpha      & N/A      & 0.5           & N/A \\
reg\_lambda     & N/A      & 2.0           & N/A \\
class\_weight   & balanced & N/A           & N/A \\
WARN threshold  & 0.50     & \textbf{0.60} & 0.65 \\
Model size      & 3.15 MB  & 1.68 MB       & 0.13 MB \\
\hline
\multicolumn{4}{l}{$^\dagger$ Converged via early stopping.} \\
\hline
\end{tabular}
\label{tab:hyper_clf}
\end{table}

%%====================================================================
%% FIGURE PLACEHOLDER A: XGBoost V2 Learning Curve
%% Sumber: notebooks/fase_8_modeling/08b_clf_xgboost.ipynb
%% Export sebagai: figures/xgb_learning_curve.png
%%====================================================================
% \begin{figure}[htbp]
%   \centering
%   \includegraphics[width=\columnwidth]{figures/xgb_learning_curve.png}
%   \caption{XGBoost V2 training and validation mlogloss vs.\ boosting
%            iteration. Convergence at iteration 499 (val\_mlogloss=0.121)
%            versus premature V1 convergence at iteration 40 (0.317).}
%   \label{fig:xgb_learning_curve}
% \end{figure}
%%====================================================================

\textbf{Model 2 --- RUL Predictor.}
Activated exclusively when Model~1 outputs WARNING or CRITICAL, justified by
four observations: (1) sensor signals in the HEALTHY zone lack discriminative
patterns for accurate RUL estimation (MAE$_\text{HEALTHY}=13.77$ vs.\
MAE$_\text{WARNING}=0.10$ days); (2) operators require short-horizon
predictions; (3) WARNING/CRITICAL detection defines the activation boundary;
and (4) conditional activation reduces LSTM frequency by $\sim$70\%.

The RUL dataset comprises 1,915 training samples (WARNING: 901, CRITICAL:
1,014), 288 validation samples, and 433 test samples (WARNING+CRITICAL only).

Three architectures were evaluated:

\begin{itemize}
\item \textit{XGBoost Regressor}: \texttt{n\_estimators=1000}
  (best\_iteration=497), \texttt{max\_depth=4}, \texttt{learning\_rate=0.01},
  \texttt{eval\_metric=mae}.

\item \textit{LSTM V2} (Selected): two-layer stacked LSTM processing
  24-timestep sequences ($\times$ 69 features). Architecture:
  LSTM(64)+L2(0.001)+Dropout(0.3)+BatchNorm $\rightarrow$
  LSTM(32)+L2(0.001)+Dropout(0.3)+BatchNorm $\rightarrow$
  Dense(16, ReLU) $\rightarrow$ Dense(1, Linear).
  Total parameters: 47,649. Best checkpoint at epoch 184 of 200.
  Optimizer: Adam (lr$=0.001$), ReduceLROnPlateau (factor$=0.5$,
  patience$=10$).

\item \textit{GRU}: GRU(48)+GRU(24)+Dropout(0.3)+BatchNorm $\rightarrow$
  Dense(16)$\rightarrow$Dense(1), $\sim$27,000 parameters. Early stopping
  at epoch 151. Overspecializes on CRITICAL (MAE$=0.014$ days) while
  performing poorly on WARNING (MAE$=1.91$ days).
\end{itemize}

%%====================================================================
%% FIGURE PLACEHOLDER D: LSTM V2 Training Learning Curve
%% Sumber: notebooks/fase_8_modeling/08e_rul_lstm.ipynb
%% Export sebagai: figures/lstm_learning_curve.png
%%====================================================================
% \begin{figure}[htbp]
%   \centering
%   \includegraphics[width=\columnwidth]{figures/lstm_learning_curve.png}
%   \caption{LSTM V2 training and validation MAE per epoch over 200
%            epochs. Best checkpoint at epoch~184 (val\_MAE$=0.8160$).}
%   \label{fig:lstm_learning_curve}
% \end{figure}
%%====================================================================

% ============================================================
\section{EXPERIMENTAL RESULTS AND DISCUSSION}
\label{sec:results}
% ============================================================

\subsection{Classification Results (Model 1)}

Tables~\ref{tab:clf_results_a}--\ref{tab:clf_results_c} compare all three
classifiers. The primary metric is \textbf{WARNING F1 on the validation set}.
Zero fatal errors (CRITICAL$\rightarrow$HEALTHY) were recorded across all
models and splits.

\begin{table}[htbp]
\caption{TABLE IXa: Classifier Accuracy Metrics}
\centering
\begin{tabular}{lcccc}
\hline
\textbf{Model} & \textbf{F1$_\text{val}$} & \textbf{F1$_\text{test}$} & \textbf{Acc$_\text{test}$} & \textbf{Fatal Err} \\
\hline
\textit{Rand.\ Forest}        & 0.9292          & 0.9914 & 0.9978 & 0 \\
\textbf{\textit{XGBoost V2}}  & \textbf{0.9894} & 0.9906 & 0.9974 & \textbf{0} \\
\textit{LightGBM}             & 0.9845          & 0.9876 & 0.9956 & 0 \\
\hline
\end{tabular}
\label{tab:clf_results_a}
\end{table}

\begin{table}[htbp]
\caption{TABLE IXb: Per-Class F1 on Test Set}
\centering
\begin{tabular}{lccc}
\hline
\textbf{Model} & \textbf{F1$_\text{HLT}$} & \textbf{F1$_\text{WARN}$} & \textbf{F1$_\text{CRIT}$} \\
\hline
\textit{Rand.\ Forest}        & 0.9998 & 0.9856 & 0.9889 \\
\textbf{\textit{XGBoost V2}}  & 0.9996 & 0.9856 & \textbf{0.9866} \\
\textit{LightGBM}             & 0.9980 & 0.9781 & 0.9865 \\
\hline
\multicolumn{4}{l}{\textbf{Selected:} XGBoost V2 (WARN F1$_\text{val}$=\textbf{0.9818})} \\
\hline
\end{tabular}
\label{tab:clf_results_b}
\end{table}

\begin{table}[htbp]
\caption{TABLE IXc: Classifier Efficiency Metrics}
\centering
\begin{tabular}{lcc}
\hline
\textbf{Model} & \textbf{Inference (ms)} & \textbf{Size (MB)} \\
\hline
\textit{Rand.\ Forest}        & 75.25          & 3.15 \\
\textbf{\textit{XGBoost V2}}  & \textbf{12.76} & 1.68 \\
\textit{LightGBM}             & 3.26           & \textbf{0.13} \\
\hline
\end{tabular}
\label{tab:clf_results_c}
\end{table}

Table~\ref{tab:threshold} shows that the default threshold of 0.50 produces
954 false alarms. Raising it to 0.60 reduces false alarms to 1 while
maintaining near-perfect WARNING recall.

\begin{table}[htbp]
\caption{TABLE X: XGBoost V2 WARNING Threshold Tuning (Validation Set)}
\centering
\begin{tabular}{ccccc}
\hline
\textbf{Threshold} & \textbf{Prec.} & \textbf{Rec.} & \textbf{F1$_\text{WARN}$} & \textbf{False Alarms} \\
\hline
0.50 (default)          & 0.087          & 0.975          & 0.159           & 954 \\
\textbf{0.60 (optimal)} & \textbf{0.973} & \textbf{0.990} & \textbf{0.9818} & \textbf{1} \\
\hline
\end{tabular}
\label{tab:threshold}
\end{table}

\subsection{RUL Predictor Results (Model 2)}

Table~\ref{tab:rul_results} presents RUL metrics as rows and model algorithms
as columns for direct cross-model comparison.

\begin{table}[htbp]
\caption{TABLE XI: RUL Predictor Comparison (WARNING+CRITICAL, $n_\text{test}$=433)}
\centering
\begin{tabular}{lp{1.5cm}p{1.5cm}p{1.5cm}}
\hline
\textbf{Metric} & \textbf{XGBoost Reg} & \textbf{LSTM V2}$^\star$ & \textbf{GRU} \\
\hline
MAE Val (days)  & 1.7189  & \textbf{1.6461}  & 1.8618 \\
MAE Test (days) & 1.1020  & \textbf{0.7985}  & 0.9515 \\
RMSE            & 5.6002  & 6.3803           & 7.3011 \\
R$^2$           & 0.3961  & 0.2594           & 0.0303 \\
E$\leq$1 day    & 93.53\% & \textbf{98.04\%} & 97.80\% \\
E$\leq$3 days   & 94.92\% & \textbf{98.04\%} & 97.80\% \\
Inference       & $\sim$116ms & $\sim$116ms & $\sim$90ms \\
\hline
\multicolumn{4}{l}{$^\star$ \textbf{Selected.} Best MAE and Error$\leq$1d.} \\
\hline
\end{tabular}
\label{tab:rul_results}
\end{table}

Table~\ref{tab:lstm_zone} reveals a zone-level asymmetry: WARNING zone
(MAE$=0.10$ days, near-perfect) vs.\ CRITICAL zone (MAE$=2.03$ days, driven
by outlier RUL values at the tail of each machine's CRITICAL window).

\begin{table}[htbp]
\caption{TABLE XII: LSTM V2 Error Analysis by Zone (Test Set)}
\centering
\begin{tabular}{lccc}
\hline
\textbf{Zone} & \textbf{N} & \textbf{MAE (days)} & \textbf{Interpretation} \\
\hline
WARNING          & 208          & \textbf{0.10}   & Near-perfect early prediction \\
CRITICAL         & 225          & 2.03            & Tail-end outlier inflation \\
\hline
\textbf{Overall} & \textbf{433} & \textbf{0.7985} & Business-grade accuracy \\
\hline
\end{tabular}
\label{tab:lstm_zone}
\end{table}

%%====================================================================
%% FIGURE PLACEHOLDER E: LSTM V2 Predicted vs Actual Scatter Plot
%% Sumber: notebooks/fase_9_evaluation/09_evaluation.ipynb
%% Export sebagai: figures/lstm_pred_vs_actual.png
%%====================================================================
% \begin{figure}[htbp]
%   \centering
%   \includegraphics[width=\columnwidth]{figures/lstm_pred_vs_actual.png}
%   \caption{LSTM V2 predicted vs.\ actual RUL ($n$=433).}
%   \label{fig:lstm_scatter}
% \end{figure}

\subsection{Discussion}

\textbf{SSBS} preserves both the prime-condition reference (Block~A) and the
pre-failure boundary (Block~B), providing contextually complete training
signal that random undersampling destroys.

\textbf{Machine-based split} ensures that every partition contains complete
failure cycles, enabling generalization to 6~unseen machines after training
on 14. Global temporal splits would concentrate all degradation examples into
the test set.

\textbf{Train-only SMOTE} guarantees that 3,288~validation and 2,716~test
records remain entirely natural, with zero influence from synthetic samples.

\textbf{LSTM vs.\ GRU}: GRU overspecializes on CRITICAL (MAE$=0.014$ days)
but fails on WARNING (MAE$=1.91$ days). LSTM's separate gates handle
WARNING-zone gradual, multi-modal degradation, justifying its 20K additional
parameters with a 19$\times$ improvement in WARNING MAE.

\textbf{R$^2$ Limitation}: CRITICAL-zone outliers (RUL up to 146.92~days)
inflate RMSE and deflate R$^2$. MAE and Error$\leq N$~days are the
business-relevant metrics; LSTM V2's lower R$^2$ (0.2594) does not indicate
inferior performance versus XGBoost (0.3961).

\textbf{Cascade efficiency}: Model~2 is invoked for only $<$13\% of ticks
(WARNING+CRITICAL), reducing total inference load by $\sim$70\% compared to
always-on sequential inference. Production smoke tests confirmed average
latency of 49.7ms (max 54.0ms), within the 200ms SLO.

\textbf{Limitations}: (1)~Synthetically generated dataset; real sensor drift
and noise may affect performance. (2)~Homogeneous fleet---heterogeneous
generalization unvalidated. (3)~Online learning not implemented.

% ============================================================
\section{CONCLUSION}
% ============================================================

PRIME addresses extreme class imbalance and temporal data leakage in a
single, unified PdM pipeline. Key contributions and results:

\begin{itemize}
\item \textbf{SSBS} reduces HEALTHY dominance from 97.36\% to 87.09\%
  while preserving time-series causality.

\item \textbf{Three-class Temporal Label Engineering} with Sensor
  Confirmation Layer (P90, minimum 2~sensors) produces validated
  HEALTHY/WARNING/CRITICAL labels.

\item \textbf{XGBoost V2} (threshold=0.60) achieves WARNING F1$=0.9818$
  with zero fatal misclassifications and 99.9\% false alarm reduction.

\item \textbf{LSTM V2} achieves MAE$=0.7985$~days and Error$\leq$1d$=98.04\%$,
  outperforming XGBoost Regressor by 27.5\% in MAE.

\item The \textbf{cascaded architecture} reduces inference overhead by
  $\sim$70\% through conditional LSTM activation.
\end{itemize}

Future work: (1)~real industrial dataset validation; (2)~online learning for
continuous adaptation; (3)~heterogeneous multi-factory generalization.

% ============================================================
% REFERENCES
% ============================================================

\begin{thebibliography}{00}

\bibitem{lee2014industryiot}
J. Lee, B. Bagheri, and H.-A. Kao, ``A Cyber-Physical Systems architecture
for Industry 4.0-based manufacturing systems,'' \textit{Manufacturing
Letters}, vol. 3, pp. 18--23, 2014.

\bibitem{zhang2019pdmsurvey}
W. Zhang, D. Yang, and H. Wang, ``Data-driven methods for predictive
maintenance of industrial equipment: A survey,'' \textit{IEEE Systems
Journal}, vol. 13, no. 3, pp. 2213--2227, Sep. 2019.

\bibitem{hochreiter1997lstm}
S. Hochreiter and J. Schmidhuber, ``Long short-term memory,''
\textit{Neural Computation}, vol. 9, no. 8, pp. 1735--1780, Nov. 1997.

\bibitem{zheng2017lstm_rul}
S. Zheng, K. Ristovski, A. Farahat, and C. Gupta, ``Long short-term memory
network for remaining useful life estimation,'' in \textit{Proc. IEEE ICPHM},
Dallas, TX, USA, 2017, pp. 88--95.

\bibitem{li2018cnnlstm}
X. Li, Q. Ding, and J.-Q. Sun, ``Remaining useful life estimation in
prognostics using deep convolution neural networks,'' \textit{Reliability
Engineering and System Safety}, vol. 172, pp. 1--11, Apr. 2018.

\bibitem{cho2014gru}
K. Cho et~al., ``Learning phrase representations using RNN encoder-decoder
for statistical machine translation,'' in \textit{Proc. EMNLP}, Doha, Qatar,
2014, pp. 1724--1734.

\bibitem{chen2016xgboost}
T. Chen and C. Guestrin, ``XGBoost: A scalable tree boosting system,'' in
\textit{Proc. 22nd ACM SIGKDD}, San Francisco, CA, USA, 2016, pp. 785--794.

\bibitem{breiman2001rf}
L. Breiman, ``Random forests,'' \textit{Machine Learning}, vol. 45, no. 1,
pp. 5--32, Oct. 2001.

\bibitem{ke2017lightgbm}
G. Ke et~al., ``LightGBM: A highly efficient gradient boosting decision
tree,'' in \textit{Advances in NeurIPS}, Long Beach, CA, USA, 2017,
pp. 3146--3154.

\bibitem{chawla2002smote}
N. V. Chawla, K. W. Bowyer, L. O. Hall, and W. P. Kegelmeyer, ``SMOTE:
Synthetic minority over-sampling technique,'' \textit{Journal of Artificial
Intelligence Research}, vol. 16, pp. 321--357, Jun. 2002.

\bibitem{li2020pdmimbalance}
Y. Li, T. Li, and J. Liu, ``Machinery fault diagnosis using deep learning
with class-imbalanced data in manufacturing,'' \textit{International Journal
of Advanced Manufacturing Technology}, vol. 109, pp. 2447--2462, 2020.

\bibitem{blagus2013smote_bias}
R. Blagus and L. Lusa, ``SMOTE for high-dimensional class-imbalanced data,''
\textit{BMC Bioinformatics}, vol. 14, p. 106, Mar. 2013.

\bibitem{bergmeir2018tscv}
C. Bergmeir, R. J. Hyndman, and B. Koo, ``A note on the validity of
cross-validation for evaluating autoregressive time series prediction,''
\textit{Computational Statistics and Data Analysis}, vol. 120, pp. 70--83,
Apr. 2018.

\bibitem{kaufman2012leakage}
S. Kaufman, S. Rosset, C. Perlich, and O. Stitelman, ``Leakage in data
mining: Formulation, detection, and avoidance,'' \textit{ACM Transactions on
Knowledge Discovery from Data}, vol. 6, no. 4, pp. 1--21, Dec. 2012.

\bibitem{hyndman2018forecasting}
R. J. Hyndman and G. Athanasopoulos, \textit{Forecasting: Principles and
Practice}, 2nd ed. OTexts, 2018. [Online]. Available:
\url{https://otexts.com/fpp2/}

\end{thebibliography}

\end{document}
